% GENERATED FILE. Edit canonical lessons, then run npm run generate:books.
% canonical-source-hash: fnv1a64:24a691ceffb07781
% canonical-lessons: ES-C387-quemar, ES-C387-cruzar, ES-C387-volar, ES-C387-girar

\chapter{Fire, Crossings, and Things That Turn}
\label{ch:quemar-y-cruzar}
\hlchaptermodality{387}

\begin{chapteropening}
\textbf{By the end of this chapter:} \emph{I can say quemar, cruzar, volar and girar, and name four everyday actions.}

Complete ES-C387-girar: name four everyday actions, and say which one you can watch happening.
\end{chapteropening}

\section[quemar]{quemar — two dictionaries, one word, no agreement}

\label{lesson:ES-C387-quemar}



\begin{quote}
Say \emph{to spend}, then \emph{to grow}. (\emph{Gastar}; \emph{crecer}.)
\end{quote}

\subsection*{What to know first}
\begin{itemize}
\raggedright
  \item el fuego — what does the burning.
  \item el humo — what it leaves behind.
\end{itemize}

\begin{sounds}
\begin{itemize}
\raggedright
  \item \emph{que-MAR}, two beats with the weight on the second. \emph{Qu} before \emph{e} is a hard \emph{k}, and the \emph{u} is silent. $\to$ pronunciation reference
\end{itemize}
\end{sounds}

\begin{cousinweb}
\textbf{quemar} — \textquotedblleft{}to burn.\textquotedblright{}

The Academy's dictionary gives this one line: from Latin \textbf{cremare}, to burn. No hedge, no \emph{quizá}. It looks closed.

Corominas opens his entry by saying it cannot be so, and his objection is entirely about sounds.

The first half of it comes from Portuguese, which has \textbf{queimar} for the same verb. That \emph{-ei-} can only have grown out of an older \textbf{ai}. \emph{Cremare} offers an \textbf{e} instead, and a short Latin \emph{e} there does something else in Spanish — it breaks into \emph{ie}. Spanish should have ended up with a \emph{quiemar}. It never has.

The second half is a missing letter. \emph{Cremare} has an \textbf{r} sitting right after the \emph{c}. \emph{Quemar} has no \emph{r} at all. Dropping one from that position is, he says, a very grave obstacle.

So he proposes that the Latin verb was reshaped along the way, under the pull of a Greek word for a burn — \emph{kaima} — carried west, he suggests, in the working vocabulary of Greek physicians, who burned wounds shut. And he says plainly that he cannot fully explain why Spanish speakers would have taken it up.

Nobody has closed this. Two dictionaries, one word, and the honest position is to hold both and pick neither.

Here is one thing about the Latin verb that is not in doubt, on the English side. English borrowed \textbf{cremation} in the sixteen-hundreds, and for two hundred years had the noun with no verb beside it. In 1851 somebody wanted the verb and, instead of borrowing one, cut the ending off the noun to make \textbf{cremate} — backwards from the usual direction of traffic.

And one to refuse. \textbf{Quench} sounds near enough to fool anybody, and means very nearly the opposite. It is Old English \emph{acwencan}, from the Germanic side, with no confident relatives outside Germanic and no Latin in it whatsoever.
\end{cousinweb}

\begin{grammarlens}[title={What you've built}]
You can now name the action: \emph{quemar}. And you can say why a phonetic objection outranks a resemblance.
\end{grammarlens}

\subsection*{Your turn}
Say these aloud:

\begin{itemize}
\raggedright
  \item \emph{quemar}
  \item \emph{quemar el papel}
  \item \emph{el fuego}
\end{itemize}

\begin{tcolorbox}[breakable,colback=teal!4,colframe=teal!35!black,title={Before you move on}]
What letter does \emph{cremare} have that \emph{quemar} does not? (An \emph{r}, right after the \emph{c}.) Is \emph{quench} a relative? (No — Old English, with no Latin in it.)
\end{tcolorbox}
\section[cruzar]{cruzar — the long way round to England}

\label{lesson:ES-C387-cruzar}



\begin{quote}
Say \emph{to grow}, then \emph{to burn}. (\emph{Crecer}; \emph{quemar}.)
\end{quote}

\subsection*{What to know first}
\begin{itemize}
\raggedright
  \item la calle — something you cross.
  \item el puente — something you cross on.
\end{itemize}

\begin{sounds}
\begin{itemize}
\raggedright
  \item \emph{cru-ZAR}, two beats with the weight on the second. The \emph{z} is a soft \emph{s} across the Americas and a \emph{th} across much of Spain. $\to$ pronunciation reference
\end{itemize}
\end{sounds}

\begin{cousinweb}
\textbf{cruzar} — \textquotedblleft{}to cross.\textquotedblright{}

It is built straight on \textbf{cruz}, and \emph{cruz} is Latin \textbf{crux}. Not on a Latin verb — Latin had one, \emph{cruciare}, and it even came to mean crossing, but the sounds do not work out, so \emph{cruzar} is a verb Spanish made for itself out of the noun it already had.

And \emph{cruz} is a slightly odd noun. Latin \emph{crux} has a short \emph{u}, and a short \emph{u} in that spot regularly comes out of Spanish as an \emph{o}. Spanish should be saying \textbf{croz}. It says \emph{cruz}, and that surviving \emph{u} is the mark of a word that stayed close to its written form because the church kept saying it — the same mark \emph{Dios} carries. A word handled by clerks and priests wears down less than one handled by everybody.

English took a very long road indeed, and the road is the story. Latin \emph{crux} went first to Ireland, where Old Irish made \textbf{cros} of it. Scandinavians picked it up from the Irish; Old English took it from them. Latin, Irish, Norse, English — three borrowings to travel a distance a Roman road already covered. The older English word for the shape was \textbf{rood}, and that is no relation.

\textbf{Crucial} is sharper than it looks. In 1620 Francis Bacon wrote of an \emph{instantia crucis}, an instance of the crossroads — the fingerpost standing where two roads fork, forcing a traveller to choose. A crucial question is a fork in a road.

\textbf{Crusade} is medieval Latin \emph{cruciata}, a thing marked with a cross; the Spanish form of that same word is \textbf{cruzada}. And \textbf{cruise} belongs here too, arriving through Dutch \emph{kruisen}, to cross, to sail to and fro. A cruise is a crossing.

One to refuse. \textbf{Cruel} is not in this family. It is Latin \emph{crudelis}, from \emph{crudus}, \emph{raw} — cruelty is rawness, not crucifixion.
\end{cousinweb}

\begin{grammarlens}[title={What you've built}]
You can now name the action: \emph{cruzar}. And you can trace one Latin word through four languages.
\end{grammarlens}

\subsection*{Your turn}
Say these aloud:

\begin{itemize}
\raggedright
  \item \emph{cruzar}
  \item \emph{cruzar la calle}
  \item \emph{el puente}
\end{itemize}

\begin{tcolorbox}[breakable,colback=teal!4,colframe=teal!35!black,title={Before you move on}]
By what road did \emph{cross} reach English? (Latin, then Old Irish, then Old Norse, then Old English.) What is a \emph{crucial} question, in Bacon's picture? (A fingerpost where two roads fork, and you have to choose.)
\end{tcolorbox}
\section[volar]{volar — a volley is a flight}

\label{lesson:ES-C387-volar}



\begin{quote}
Say \emph{to burn}, then \emph{to cross}. (\emph{Quemar}; \emph{cruzar}.)
\end{quote}

\subsection*{What to know first}
\begin{itemize}
\raggedright
  \item el avión — something that flies.
  \item el cielo — where it flies.
\end{itemize}

\begin{sounds}
\begin{itemize}
\raggedright
  \item \emph{vo-LAR}, two beats with the weight on the second. The \emph{v} is the same sound as \emph{b}. $\to$ pronunciation reference
\end{itemize}
\end{sounds}

\begin{cousinweb}
\textbf{volar} — \textquotedblleft{}to fly.\textquotedblright{}

Latin \textbf{volare}, to fly, and nothing has been bent on the way. So the payoff is in two English words that stopped sounding like flight.

A \textbf{volley} is the better one. It came through a French word meaning \emph{a flight} — and the first volleys were arrows. A volley of arrows is not a metaphor and never was: it is a \emph{flight} of them, released together. Centuries later a tennis player borrowed the word for a ball struck while still in the air, which is the same picture with one object instead of a hundred.

\textbf{Volatile} is the other. Latin \emph{volatilis} meant \emph{flying, winged}. English now uses it for a liquid that disappears into the air and for a person whose temper does the same. Both are the Latin word doing its ordinary job: something that gets away from you.

One caution, and it is a Spanish one rather than an English one. \emph{Volar} and \emph{volver} look like each other and are not related. They descend from two different Latin verbs that happen to have arrived at three shared letters. The resemblance is an accident, and treating it as a family tie is the exact mistake this strand exists to stop.
\end{cousinweb}

\begin{grammarlens}[title={What you've built}]
You can now name the action: \emph{volar}. And you know what the first volleys were.
\end{grammarlens}

\subsection*{Your turn}
Say these aloud:

\begin{itemize}
\raggedright
  \item \emph{volar}
  \item \emph{el avión}
  \item \emph{cruzar el cielo}
\end{itemize}

\begin{tcolorbox}[breakable,colback=teal!4,colframe=teal!35!black,title={Before you move on}]
What were the first volleys? (Arrows — a flight of them.) What did Latin \emph{volatilis} mean? (Flying, winged.)
\end{tcolorbox}
\section[girar]{girar — a word that came in through a book}

\label{lesson:ES-C387-girar}



\begin{quote}
Say \emph{to cross}, then \emph{to fly}. (\emph{Cruzar}; \emph{volar}.)
\end{quote}

\subsection*{What to know first}
\begin{itemize}
\raggedright
  \item la llave — something you turn.
  \item la derecha — a direction you turn in.
\end{itemize}

\begin{sounds}
\begin{itemize}
\raggedright
  \item \emph{gi-RAR}, two beats with the weight on the second. \emph{G} before \emph{i} is the harsh throat sound, the same one \emph{j} makes. $\to$ pronunciation reference
\end{itemize}
\end{sounds}

\begin{cousinweb}
\textbf{girar} — \textquotedblleft{}to turn.\textquotedblright{}

The trail runs Latin \textbf{gyrare}, from \textbf{gyrus}, a ring or a circle — and Latin had that from Greek \textbf{gyros}, the same thing.

But \emph{girar} is unusual company for the other verbs around it, and the reason is worth knowing. Every one of them came down through ordinary speech, mouth to mouth, wearing away as it went. \emph{Girar} did not. Corominas calls it frankly bookish: scholars lifted it out of written Latin and put it into Spanish. It turns up in the middle of the fourteen-hundreds among writers straining to sound Latin, and then almost vanishes for a century and a half. As late as 1611 Covarrubias was still telling readers it was \textbf{not a word used in Castile}.

An inherited word arrives worn. A borrowed one arrives clean. \emph{Girar} looks almost exactly like \emph{gyrare} because it never had the centuries of use that would have rubbed it down.

The Greek circle is still busy in English. A \textbf{gyroscope} is a circle-watcher — \emph{gyros} plus \emph{skopos}, a watcher — named in French in 1852 for a device whose whole point is watching a spin.

And the \textbf{gyros} you can buy for lunch is the same Greek word, meaning \emph{turn} — named for meat turning on an upright spit. It is a genuine relative, but not by this road: it walked into English straight from modern Greek in 1971 and never went through Latin at all. The Greek name is itself a translation of the Turkish name for the same food, which also means \emph{turning}. Three languages, one idea, and none of them borrowed the meat without borrowing the spin.
\end{cousinweb}

\begin{grammarlens}[title={What you've built}]
You can now name the action: \emph{girar}. And you can tell an inherited word from a borrowed one by how worn it looks.
\end{grammarlens}

\subsection*{Your turn}
Say these aloud:

\begin{itemize}
\raggedright
  \item \emph{girar}
  \item \emph{girar la llave}
  \item \emph{la derecha}
\end{itemize}

\begin{tcolorbox}[breakable,colback=teal!4,colframe=teal!35!black,title={Before you move on}]
Did \emph{girar} come down through speech or out of a book? (Out of a book — it was still being called unusual in 1611.) Did the sandwich come through Latin? (No — straight from modern Greek, and only in 1971.)
\end{tcolorbox}
